\section{半平面交}

\subsection{半平面的存储}
\begin{frame}
\pause
一条直线把平面划分成两个区域，其中的任意一个区域就叫做半平面。

\pause
半平面一般用带顺序的两点式直线来表示。对于直线(a,b)，我们默认它代表的半平面为向量ab沿逆时针方向扫描180$^{o}$经过的区域。
\end{frame}

\subsection{半平面交}
\begin{frame}
\pause
即给定若干个半平面，求出它们的公共部分。

\pause
求解半平面交的方法是：首先把所有表示半平面的直线向量进行极角排序，极角相同的把被包含的那个保留，其它的删除。

\pause
然后按照极角从小往大加入半平面。并用一个deque维护当前已经加入的半平面。加入一个半平面时，如果队列尾部和尾部前一个元素的交点不在当前半平面内，则弹出队列尾。同时也检查队首，如果队列首部和首部下一个元素不在当前半平面内也将其删除。

\pause
在加入完所有的半平面后，还要做一次队列首尾交接的工作。如果队列首部和首部的下一个元素不在队尾的半平面内则把队首弹出。如果队列尾部和尾部的前一个元素不在队首的半平面内则把队尾弹出
\end{frame}